\chapter{Seiten und Navigation}\label{kap:seiten}

\section{Grundprinzip: statische Seiten}

Die Website besteht aus statischen, in sich geschlossenen HTML-Seiten. Der
Scheduler schreibt sie nach jedem Simulationszyklus neu, ein Webserver
(Caddy) liefert sie aus. Es gibt keinen Anwendungsserver und keine Datenbank.
Bis August 2026 lief die Seite als Shiny-App auf shinyapps.io. ADR~0001
begründet den Wechsel: Der Dienst endet im März 2027, die App hatte keine
echte Interaktivität, die Server-Rechenzeit gebraucht hätte, und an einem
Spieltag waren bis zu 240 Neuveröffentlichungen nötig.

Die Heatmap ist eine HTML-Tabelle mit eingefärbten Zellen, es gibt keine
Bilder. Alle Seiten zusammen sind einige hundert Kilobyte groß. Nur das
ELO-Verlaufsdiagramm zeichnet ein Skript im Browser.

\section{Seiteninventar}

Pro Zyklus entstehen 22 Seiten
(\datei{RCode/generate_static_site.R:1086--1216}): zehn Ligaseiten, zehn
ELO-Verlaufsseiten, die Aufstiegsseite der Regionalligen und die
Methodik-Seite. Liegen gar keine Ergebnisse vor, entsteht nur eine
Platzhalter-Startseite (\enquote{Noch keine Prognosedaten verfügbar}).

\begin{center}
\footnotesize
\setlength{\tabcolsep}{4pt}
\begin{tabular}{@{}lllr@{}}
\toprule
Liga & Ligaseite & Verlaufsseite & ID \\
\midrule
Bundesliga & \code{index.html} & \code{bundesliga-verlauf.html} & 78 \\
2.~Bundesliga & \code{2-bundesliga.html} & \code{2-bundesliga-verlauf.html} & 79 \\
3.~Liga & \code{3-liga.html} & \code{3-liga-verlauf.html} & 80 \\
RL Nord & \code{rl-nord.html} & \code{rl-nord-verlauf.html} & 84 \\
RL Nordost & \code{rl-nordost.html} & \code{rl-nordost-verlauf.html} & 85 \\
RL West & \code{rl-west.html} & \code{rl-west-verlauf.html} & 87 \\
RL SüdWest & \code{rl-suedwest.html} & \code{rl-suedwest-verlauf.html} & 86 \\
RL Bayern & \code{rl-bayern.html} & \code{rl-bayern-verlauf.html} & 83 \\
Frauen-BL & \code{frauen-bundesliga.html} & \code{frauen-bundesliga-verlauf.html} & 82 \\
2.~Frauen-BL & \code{2-frauen-bundesliga.html} & \code{2-frauen-bundesliga-verlauf.html} & 1034 \\
\midrule
Aufstieg & \code{rl-aufstieg.html} & -- & -- \\
Methodik & \code{methodik.html} & -- & -- \\
\bottomrule
\end{tabular}
\end{center}

Die Bundesliga ist die Startseite und heißt deshalb \code{index.html}. Ihre
Verlaufsseite heißt \code{bundesliga-verlauf.html} und nicht
\code{index-verlauf.html}. Alle Namen stammen aus der Liga-Registry
(\datei{RCode/league_registry.R}) und aus \code{league\_views()}
(\datei{RCode/league_views.R}). Die Registry hält alles Ligaspezifische als
Daten (ADR~0005). Eine neue Liga ist dort ein neuer Eintrag und keine neue
Code-Verzweigung.

\section{Zweistufige Navigation}

Jede Seite trägt denselben Kopf: die Wortmarke \enquote{30 Punkte}, darunter
einen kurzen roten Strich, die Unterzeile und die Navigation. Die Navigation
hat eine Zeile je Gruppe. Jede Zeile beginnt mit dem Gruppennamen in
Kapitälchen, danach folgen die Links, getrennt durch Mittelpunkte:

\begin{center}
\begin{tabular}{@{}>{\sffamily\scriptsize\color{ink3}}l>{\sffamily\small}l@{}}
HERREN       & Bundesliga · 2.~Bundesliga · 3.~Liga \\
REGIONALLIGA & Nord · Nordost · West · SüdWest · Bayern · Aufstieg \\
FRAUEN       & Bundesliga · 2.~Bundesliga \\
             & Methodik \\
\end{tabular}
\end{center}

Die Regionalligen stehen bewusst direkt unter den Herren
(\datei{generate_static_site.R:54--69}, Issue~\#178). Sie sind
Herren-Ligen derselben Wechselgemeinschaft, und die Nachbarschaft soll das
zeigen. Die aktuelle Seite ist fett und unterstrichen und trägt
\code{aria-current="page"}.

Auf den Verlaufsseiten führen die Liga-Links zur \emph{Verlaufsseite} der
jeweils anderen Liga. Wer Verläufe vergleicht, bleibt so in dieser Ansicht.
Aufstieg und Methodik verlinken normal.

\section{Hinweis auf veraltete Daten}

Eine statische Seite veraltet genau dann, wenn der Scheduler nicht mehr
schreibt. Eine Prüfung beim Rendern würde deshalb nie anschlagen. Die
Entscheidung fällt im Browser (\datei{generate_static_site.R:445--469}):
Im Fuß jeder Seite steht der Zeitpunkt der Erzeugung als
\code{<time id="generated" datetime="…">} in UTC. Ein kleines Skript
berechnet daraus das Alter in Stunden,
\[
  \text{Alter} = \frac{\text{jetzt} - \text{erzeugt}}{3{,}6\cdot 10^{6}\ \text{ms}},
\]
und blendet ab einem Alter von mehr als 24 Stunden einen Hinweis ein:
\enquote{Achtung: Diese Prognosen sind \emph{n} Stunden alt und werden derzeit
nicht aktualisiert.}

Daneben gibt es einen zweiten, serverseitigen Hinweis. Arbeitet der
Scheduler wegen knappen API-Kontingents in einem gedrosselten Takt (mehr als
180 Sekunden statt 120), ergänzt der Fuß \enquote{Eingeschränkter Service,
Update etwa alle \emph{n} Minuten}. Anlass war ein Vorfall am 14.~September
2026, bei dem die Seite nur alle 83 Minuten aktualisiert wurde.
